\section{Tecnologías y Herramientas Utilizadas}

La elección de las tecnologías empleadas responde a la necesidad de crear un sistema híbrido que combine la robustez transaccional con la agilidad en el análisis de datos \parencite{hilpisch_python_2018}. A continuación, se detallan las herramientas principales y su justificación técnica:

\begin{description}
    \item[Java 21 y Spring Boot 3:] Se ha seleccionado como lenguaje de \textit{orquestador Java} principal ebido a su fuerte tipado y su modelo de concurrencia multihilo, esencial para gestionar múltiples strategias de trading simultáneas sin errores de memoria \parencite{treleaven_algorithmic_2013}. El cosistema Spring Boot 3 facilita la creación de una arquitectura modular mediante la inyección de ependencias y simplifica la persistencia de datos con Spring Data JPA \parencite{pivotal_springboot_2024}. La adopción de Java 21 permite el uso de hilos virtuales (Project Loom) \parencite{goetz_jep444_2023}, una innovación de la plataforma que proporciona millones de hilos ligeros sin la sobrecarga de los hilos del sistema operativo tradicionales.
    \item[Python 3.11+:] Es el estándar de facto en la industria para la ciencia de datos y el aprendizaje automático \parencite{goodfellow_deep_2016}. Su sintaxis flexible permite iterar rápidamente en el diseño de estrategias lógicas \parencite{hilpisch_python_2018}. En este proyecto, Python actúa como el motor de cálculo matemático, delegando en él las tareas de procesamiento de señales, descarga de datos de mercado y entrenamiento de modelos \parencite{mckinney_python_2022}.
    \item[Apache Maven:] Utilizado como gestor de dependencias y herramienta de construcción para el componente Java. Permite estandarizar el ciclo de vida del proyecto, garantizando que todas las librerías se gestionen de forma centralizada y reproducible \parencite{apache_maven_2024}.
    \item[MySQL 8.0+:] Como sistema de gestión de bases de datos relacionales (RDBMS). Su elección se justifica por su madurez y su excelente soporte para transacciones ACID \parencite{oracle_mysql_acid_2024}, lo cual es crítico para asegurar que el libro mayor y los saldos de las carteras digitales nunca queden en un estado inconsistente tras una operación técnica \parencite{oracle_mysql_2024}.
    \item[MessagePack:] Formato binario de serialización utilizado para la comunicación inter-proceso (IPC) entre Java y Python. A diferencia de JSON, MessagePack ofrece un tamaño de mensaje sustancialmente menor y permite el uso de prefijado con longitud explícita, garantizando que los límites de los mensajes sean explícitos y no se corra el riesgo de lecturas incompletas o solapadas en tuberías de comunicación \parencite{furuhashi_msgpack_2013}.
\end{description}

\subsection{Librerías Críticas}

Para alcanzar los objetivos del proyecto, se han integrado librerías especializadas que extienden las capacidades de los lenguajes base \parencite{richards_software_2022}.

\subsubsection{Stack Java}

\begin{itemize}
    \item \textbf{Spring Data JPA:} Proporciona una capa de abstracción sobre Hibernate, permitiendo el uso del patrón Repository \parencite{gamma_design_1994} para consultas y persistencia sin escribir SQL manual en la mayoría de casos \parencite{spring_data_jpa_2024}.
    \item \textbf{Lombok:} Utilizada para reducir el código repetitivo en las entidades y objetos de transferencia de datos, mejorando la legibilidad y el mantenimiento del código fuente mediante la generación automática de constructores por inyección de dependencias y registro estructurado de eventos \parencite{richards_software_2022,lombok_project_2024}.
    \item \textbf{MessagePack for Jackson:} Integración entre Jackson (el serializador por defecto de Spring) y el formato MessagePack \parencite{furuhashi_msgpack_2013}, permitiendo serializar y deserializar el protocolo IPC de forma eficiente \parencite{pivotal_springboot_2024}.
    \item \textbf{Dotenv-java:} Permite la gestión segura de variables de entorno, separando la configuración sensible (credenciales de base de datos o claves de API) del código fuente \parencite{owasp_password_2024}.
    \item \textbf{CLI interactiva:} La interfaz de línea de comandos interactiva se implementa directamente sobre las primitivas de arranque de Spring Boot 3 \parencite{pivotal_springboot_2024}, ejecutando un bucle de lectura de entrada estándar que parsea comandos en texto plano. Esta aproximación prescinde de \textit{frameworks} CLI adicionales, reduciendo dependencias externas y manteniendo el control total sobre el ciclo de vida de la sesión interactiva. 
    \item \textbf{HikariCP:} Pool de conexiones a base de datos de alto rendimiento, optimizado para aplicaciones de baja latencia \parencite{wooldridge_hikaricp_2024}.
\end{itemize}

\subsubsection{Stack Python}

\begin{itemize}
    \item \textbf{Pandas y NumPy:} Fundamentales para la manipulación eficiente de las series temporales de precios \parencite{mckinney_python_2022}. Permiten realizar operaciones vectorizadas sobre miles de velas en milisegundos, algo vital para la simulación histórica de estrategias \parencite{treleaven_algorithmic_2013}.
    \item \textbf{Scikit-Learn, XGBoost y LightGBM:} Proporcionan las implementaciones del estado del arte para los modelos predictivos \parencite{pedregosa_scikit_2011}. XGBoost, en particular, se utiliza por su capacidad de regularización para filtrar indicadores ruidosos \parencite{chen_xgboost_2016}. LightGBM complementa el catálogo permitiendo entrenamientos más rápidos en conjuntos de datos de gran volumen \parencite{ke_lightgbm_2017}.
    \item \textbf{TensorFlow/Keras:} Plataforma de aprendizaje profundo utilizada para la construcción y entrenamiento de redes neuronales densas en los modelos de predicción de series temporales \parencite{goodfellow_deep_2016}.
    \item \textbf{Optuna:} Motor de optimización bayesiana de hiperparámetros \parencite{feurer_hyperparameter_2019}. A partir de un espacio de búsqueda parametrizable, explora de forma guiada las configuraciones más prometedoras y descarta automáticamente los ensayos de baja probabilidad de éxito, reduciendo el tiempo de búsqueda frente a la exploración exhaustiva. Se utiliza en el motor de optimización del proyecto para encontrar la configuración óptima de cada modelo predictivo.
    \item \textbf{Joblib:} Biblioteca de serialización utilizada para almacenar en disco los modelos entrenados. Permite guardar y recuperar el estado completo de un modelo con una sola llamada, lo que hace posible persistir los resultados de cada sesión de entrenamiento y reutilizarlos directamente en la predicción en tiempo real sin necesidad de reentrenar.
    \item \textbf{Requests:} Librería de cliente HTTP utilizada para consultar la interfaz de programación de Binance al descargar datos históricos de mercado \parencite{binance_rest_api_2024,requests_python_2024}.
    \item \textbf{MessagePack:} Librería de serialización correspondiente a Python, permitiendo la comunicación simétrica con el lado Java \parencite{furuhashi_msgpack_2013}.
    \item \textbf{TA:} Proporciona implementaciones de indicadores técnicos clásicos como el RSI, las medias móviles, el MACD y el ATR \parencite{murphy_technical_1999}, integrados directamente sobre las tablas de datos de Pandas \parencite{mckinney_python_2022} sin dependencias nativas de compilación.

\end{itemize}

\subsubsection{Herramientas de Calidad y Asistencia al Desarrollo}

\begin{itemize}
    \item \textbf{GitHub Copilot:} Asistente de programación basado en
    inteligencia artificial desarrollado por GitHub (Microsoft)
    \parencite{github_copilot_2024}. Durante el desarrollo del proyecto se ha
    utilizado para la generación asistida de código repetitivo, la propuesta
    de pruebas unitarias y la exploración de variantes de implementación, siempre
    bajo supervisión del desarrollador para garantizar la corrección semántica
    y la coherencia con la arquitectura hexagonal establecida.

    \item \textbf{SonarLint:} Herramienta de análisis estático integrada en el entorno de desarrollo \parencite{sonarsource_sonarlint_2024}. Proporciona retroalimentación en tiempo real sobre problemas de código, vulnerabilidades y violaciones de buenas prácticas Java directamente en el editor, adelantando la detección de problemas al momento de escritura del código antes de la fase de construcción del proyecto.

    \item \textbf{Checkstyle, PMD y SpotBugs:} Conjunto de herramientas de
    análisis estático integradas en el ciclo de construcción Maven.
    Checkstyle (v10.17.0) verifica la consistencia del estilo de codificación;
    PMD (v3.24.0) detecta fragmentos de código potencialmente problemáticos
    mediante análisis del árbol sintáctico abstracto; SpotBugs (v4.8.6)
    identifica patrones de errores conocidos mediante análisis del código compilado \parencite{owasp_password_2024}. Estas herramientas se configuran como barreras de calidad en la fase de verificación de Maven, impidiendo que código con violaciones supere la construcción.

    \item \textbf{JaCoCo:} Extensión Maven de medición de cobertura de pruebas
    (v0.8.12) que genera informes de cobertura por rama e instrucción tras
    cada ejecución de las pruebas automatizadas, permitiendo identificar áreas del sistema no ejercitadas.

\end{itemize}